\subsection{Both edge directions and reference edge data}
\label{subsec:peps_arc_window_gauge}

\begin{definition}[Exchange of torus coordinates]%
    \label{def:peps_torus_coordinate_swap}
    \lean{TNLean.PEPS.torusCoordinateSwap}
    \lean{TNLean.PEPS.torusCoordinateSwap_apply}
    \lean{TNLean.PEPS.torusCoordinateSwap_symm}
    \leanok
    The coordinate exchange $S(x,y)=(y,x)$ is a graph isomorphism from
    the $w\times h$ torus to the $h\times w$ torus. Its inverse is the
    coordinate exchange in the opposite direction.
\end{definition}

\begin{lemma}[Arc windows and translations under coordinate exchange]%
    \label{thm:peps_torus_coordinate_swap_windows}
    \lean{TNLean.PEPS.Region_map_torusCoordinateSwap_arcRectangle}
    \lean{TNLean.PEPS.NormalTorusArcWindowInjectivityHypotheses.transportCoordinateSwap}
    \lean{TNLean.PEPS.torusCoordinateSwap_translate}
    \lean{TNLean.PEPS.Edge.map_torusCoordinateSwap_rightEdge}
    \lean{TNLean.PEPS.IsTorusTranslationInvariant.transportCoordinateSwap}
    \leanok
    Coordinate exchange sends the cyclic $L\times K$ rectangle starting
    at $(x,y)$ to the cyclic $K\times L$ rectangle starting at $(y,x)$,
    and sends the right edge at $(x,y)$ to the up edge at $(y,x)$.
    It intertwines translation by $(a,b)$ with translation by $(b,a)$.
    Consequently, translation invariance and injectivity of every
    cyclic $L\times K$ window are preserved, with the side lengths
    exchanged.
\end{lemma}
\begin{proof}
    \leanok
    \uses{def:peps_torus_coordinate_swap, thm:peps_graph_transport_composition,
        thm:peps_region_injective_identification}
    Membership in the two coordinate arcs is exchanged, and relabelling
    along a graph isomorphism preserves regional injectivity. The identity
    $S\circ t_{a,b}=t_{b,a}\circ S$ and functoriality of graph relabelling
    give translation invariance of the relabelled tensor.
\end{proof}

\begin{lemma}[Equality of all bond dimensions from normal arc windows]%
    \label{thm:peps_arc_window_all_bond_dimensions}
    \lean{TNLean.PEPS.bondDim_eq_of_normalArcWindows_sameState}
    \leanok
    \uses{def:peps_normal_torus_arc_window_tensor}
    Let $A$ and $B$ be translation-invariant PEPS tensors on the
    $w\times h$ torus with the same contracted state, both normal at
    window size $L\times K$
    (Definition~\ref{def:peps_normal_torus_arc_window_tensor}), where
    $L,K>0$, $2L+1\leq w$, and $2K+1\leq h$. Then $D_A(e)=D_B(e)$ at
    every edge $e$.
\end{lemma}
\begin{proof}
    \leanok
    \uses{thm:peps_staircase_cross_tensor_bond_dimension,
        thm:peps_torus_coordinate_swap_windows,
        thm:peps_torusUniformBondDim_of_translationInvariant}
    Without assuming any equality of non-boundary bond products,
    applying the staircase bond-dimension comparison to a horizontal
    reference edge and again after coordinate exchange (with side lengths
    $K,L$) establishes equality at both a horizontal and a vertical reference
    edge. Since translation invariance makes the bond dimension constant
    within each orientation, these two comparisons determine every
    edge dimension.
\end{proof}

\begin{lemma}[A covariant absorbed gauge family from two reference edges]%
    \label{thm:peps_absorbed_family_reference_witnesses}
    \lean{TNLean.PEPS.exists_torusCovariantAbsorbedGaugeFamily_of_edgeReferenceWitnesses}
    \leanok
    Let $A,B$ be translation-invariant tensors on a torus with
    $w,h>2$, matched bond dimensions, and positive bond dimensions
    for $A$. Let $e_{\mathrm h}$ and $e_{\mathrm v}$ be distinguished
    horizontal and vertical edges. Suppose that at each of these
    edges there are reference coefficient-identity data: a region $R$ (not
    necessarily a rectangle) such that
    $R$ and its complement are injective for $B$, all bond dimensions of $B$ are
    positive, and two invertible matrices $Z,Z_{\mathrm{ref}}$ satisfy
    $C_A(R,e,M;\sigma,\tau)=C_B(R,e,Z M Z^{-1};\sigma,\tau)$
    for both choices $Z$ and $Z_{\mathrm{ref}}$, all matrices, and all
    physical configurations, after identifying
    the matched matrix sizes. Require that the first stored endpoint
    of each distinguished edge belongs to its region $R$.
    Then there is a translation-covariant family of invertible edge
    matrices $X_e$ such that, for every edge $e$, global physical
    configuration $\sigma$, and matrix $M$,
    \begin{align}
        C_A(e,M;\sigma)=C_{X\cdot B}(e,M;\sigma),
    \end{align}
    where $X\cdot B$ denotes the tensor obtained by the edge gauges
    and the matrix on the right is identified by the matched bond
    dimensions.
\end{lemma}
\begin{proof}
    \leanok
    \uses{def:peps_absorbedBoundaryGauge,
        thm:peps_regionInsertedCoeff_translate_coeffIdentity}
    Translating the horizontal reference data to each horizontal edge and
    the vertical reference data to each vertical edge allows the gauge on
    each translated boundary to be absorbed into the tensor $B$.
    Because the translation reaching an edge is unique when $w,h>2$,
    composition of translations gives covariance of the
    resulting family, with the absorbed matrix replaced by its transposed
    inverse whenever the stored endpoint order is reversed.
    The translated coefficient identity then gives the displayed
    bare-edge identity at each edge.
\end{proof}
